\documentclass[10pt,a4paper]{amsart}
\usepackage[margin=19mm]{geometry}
\usepackage{amsmath,amssymb,mathtools}
\usepackage[scr=rsfs,cal=euler]{mathalfa}
\usepackage{xfrac}
\usepackage[unicode,hidelinks,bookmarks=false]{hyperref}
\newcommand{\ie}{i.e.\ }
\newcommand{\cf}{cf.\ }
\newcommand{\resp}{respectively\ }
\newcommand{\Subsection}{\subsection}
\providecommand{\nobreakdash}{\nobreak\hbox{-}}
\setlength{\emergencystretch}{2em}
\multlinegap0pt
\usepackage{amssymb}
\usepackage{xcolor}
\usepackage{tikz}
\newtheorem{proposition}{Proposition}
		\newcommand{\Ab}{{\mathbf{Ab}}} %abelian groups
\newcommand{\BI}{{\operatorname{BI}}} % Bredon-Illman
\newcommand{\Coeff}{{\operatorname{Coeff}}} % coefficient systems
		\newcommand{\calA}{{\mathcal{A}}}
		\newcommand{\calB}{{\mathcal{B}}}
\newcommand{\labell}[1]{\label{#1} \printname{#1}}
\newcommand{\printname}[1]{}
\title{Twisted Bredon-Illman Cohomology is a Morita Invariant}
\author{Carla Farsi, Laura Scull,, Jordan Watts}
\date{Source: arXiv:2507.03091v2 (CC BY 4.0)}
\begin{document}
\maketitle

\noindent\emph{Source and licence.} Twisted Bredon-Illman Cohomology is a Morita Invariant. Version arXiv:2507.03091v2 (CC BY 4.0), distributed by arXiv under the Creative Commons Attribution 4.0 International licence, http://creativecommons.org/licenses/by/4.0/, as recorded in the arXiv licence metadata for this exact version and shown on the abstract page https://arxiv.org/abs/2507.03091v2. Full licence notice: \texttt{licenses/arxiv-2507.03091-CC-BY-4.0.txt}.\par\smallskip

\noindent\textit{Editorial scope.} Counted unit: the article's proposition proposition (source0.tex lines 1369--1371). The article's complete original proof of it is delivered with it (source0.tex lines 1373--1385, one \texttt{proof} environment of 13 lines). No mathematics is written, repaired or re-derived: the statement and the proof are the article's own lines, in the article's own notation, and no retained line is substituted or re-flowed.  No further source-local context is needed: every cross-reference of every retained line resolves inside the two delivered spans.  Nothing was omitted from any retained span, so the package carries no omission record.  No retained line contains a citation, so the wrapper carries no bibliography and no bibliography entry.  No retained line contains a figure environment, an image-inclusion command or drawing code, so no image is delivered and the package declares no asset dependency.  Editorial formatting only, in addition: an AMS \texttt{amsart} preamble that restates the article's own declarations and macros used by the retained lines, namely the environments \texttt{proposition} on separate counters, together with the standard packages those lines need.  The retained environments print in this document as Proposition 1 in order of appearance, whereas the article prints the same items with the numbering of its own full document.  The complete native source of the article as distributed in the arXiv e-print for this exact version is bundled unchanged as \texttt{sources/180993/source0.tex}.
    \begin{proposition}[Functoriality of $H_\BI^\bullet$ in Coefficient Systems]\labell{p:nat isom cohom}
        For any $G$-space $X$ and coefficient systems $\calA$ and $\calB$, a natural transformation $\eta\colon\calA\Rightarrow\calB$ induces a map of cohomology groups $\eta_*\colon H^\bullet_\BI(G\ltimes X;\calA)\to H^\bullet_\BI(G\ltimes X;\calB)$.  This makes $H^\bullet_\BI(G\ltimes X;\cdot)$ into a functor from $\Coeff(G\ltimes X)$ to $\Ab$.  In particular, if $\eta$ is a natural isomorphism, then $\eta_*$ is an isomorphism of cohomology groups.
    \end{proposition}

    \begin{proof}
        Suppose $\eta\colon\calA\Rightarrow\calB$ is a natural isomorphism.  Define $\eta_*\colon S_\BI^n(G\ltimes X;\calA)\to C_\BI^n(G\ltimes X;\calB)$ as follows: for $\sigma\colon\Delta_n\times G/K\to X$ an equivariant $n$-simplex, let $$\eta_*(c)(\sigma):=\eta_{\sigma_K}(c(\sigma)).$$  It follows from the naturality of $\eta$ that $\eta_*c\in S_\BI^n(G\ltimes X;\calB)$. 

        Fix $\sigma\colon\Delta_{n+1}\times G/K\to X$ an equivariant $(n+1)$-simplex, and $c\in S_\BI^n(G\ltimes X;\calA)$.
        \begin{align*}
            \eta_*(\delta c)(\sigma)  =&~ \eta_*\left(\calA([\sigma_K^{(0)}])c(\sigma^{(0)})+\sum_{j=1}^{n+1}(-1)^jc(\sigma^{(j)})\right) \\
                =&~ \calB([\sigma_K^{(0)}])\eta_*c(\sigma^{(0)})+\sum_{j=1}^{n+1}(-1)^j\eta_*c(\sigma^{(j)}) \\
                =&~ \delta(\eta_*c)(\sigma)
        \end{align*}
        It follows that $\eta$ induces a homomorphism $\eta_*\colon H_\BI^\bullet(G\ltimes X;\calA)\to H_\BI^\bullet(G\ltimes X;\calB).$  

        The assignment $\calA\to H^\bullet_\BI(G\ltimes X,\calA)$ and $(\eta\colon\calA\to\calB)\mapsto\eta_*$ is functorial; in particular, if $\eta$ is a natural isomorphism, then $\eta_*$ is an isomorphism (in fact, on the level of cochains).  The result follows.
    \end{proof}

\end{document}
